\section{Sistemas de primer orden $u' = Au + B(t)$}

Para el estudio de sistemas de ecuaciones diferenciales ordinarias, donde $A \in \mathcal{M}_{n \times n}(\mathbb{C})$ es una matriz con coeficientes constantes, la teoría se desarrolla de manera análoga buscando soluciones para el sistema homogéneo asociado:
\[
    h' = Ah
\]

Aprovechando la estructura lineal, proponemos una solución vectorial de la forma $h(t) = e^{\lambda t}v$, donde $\lambda \in \mathbb{C}$ es un escalar y $v \in \mathbb{C}^n \setminus \{\vec{0}\}$ es un vector constante.
Derivando e introduciendo la propuesta en el sistema, obtenemos:
\[
    \lambda e^{\lambda t} v = A (e^{\lambda t} v) \implies e^{\lambda t} (\lambda v) = e^{\lambda t} (Av)
\]
Como $e^{\lambda t} \neq 0$, simplificamos para llegar a la ecuación fundamental:
\[
    Av = \lambda v
\]

Esta ecuación nos indica que $\lambda$ debe ser un \textbf{autovalor} de la matriz $A$ y $v$ debe ser su respectivo \textbf{autovector} asociado. Al conjunto de todos los autovalores de la matriz $A$ lo denotamos como $\sigma(A)$ (el \textit{espectro} de $A$).
\[
    \sigma(A) = \{\lambda_1, \lambda_2, \dots, \lambda_k\}
\]
La condición necesaria y suficiente para que exista un vector no nulo $v$ que satisfaga $Av = \lambda v$ es que la matriz $(A - \lambda I)$ no sea invertible, lo cual ocurre si y solo si su determinante se anula. Esto nos lleva a la ecuación característica de la matriz:
\[
    \lambda \in \sigma(A) \iff \det(A - \lambda I) = 0
\]
Desarrollando este determinante, obtenemos un polinomio de grado $n$ en $\lambda$, que guarda una estrecha relación con el caso escalar unidimensional:
\[
    \det(A - \lambda I) = (-1)^n \left( \lambda^n - a_{n-1} \lambda^{n-1} - \dots - a_1 \lambda - a_0 \right) = 0
\]
Una base del espacio de soluciones del sistema homogéneo $h' = Ah$ se construye a partir de los autovectores asociados a cada autovalor. Si $\sigma(A) = \{\lambda_1, \lambda_2, \dots, \lambda_n\}$ con $\lambda_j \neq \lambda_k$ para $j \neq k$, y $v_j$ es un autovector asociado a $\lambda_j$, entonces el conjunto $\{e^{\lambda_1 t} v_1, e^{\lambda_2 t} v_2, \dots, e^{\lambda_n t} v_n\}$ forma una base de soluciones del sistema homogéneo.
\begin{lema}
    Sean $\lambda_1, \lambda_2, \dots, \lambda_n$ números complejos distintos dos a dos, y sean $c_1, c_2, \dots, c_n$ constantes tales que 
    \[
        c_1 e^{\lambda_1 t} + c_2 e^{\lambda_2 t} + \dots + c_n e^{\lambda_n t} = 0 \quad \forall t \in \mathbb{R}.
    \]
    Entonces, necesariamente $c_j = 0$ para todo $j \in \{1, 2, \dots, n\}$.
\end{lema}

\begin{proof}
    Demostraremos el resultado por inducción sobre $n$.
    
    \textbf{Caso base ($n=1$):} La ecuación se reduce a $c_1 e^{\lambda_1 t} = 0$. Dado que la función exponencial nunca se anula ($e^{\lambda_1 t} \neq 0$ para todo $t$), se concluye directamente que $c_1 = 0$.
    
    \textbf{Hipótesis de inducción:} Supongamos que el resultado es cierto para $n-1$ términos. Es decir, si $c_1 e^{\lambda_1 t} + \dots + c_{n-1} e^{\lambda_{n-1} t} = 0$ para todo $t$, entonces $c_j = 0$ para todo $j \in \{1, 2, \dots, n-1\}$.
    
    \textbf{Paso inductivo:} Consideremos la ecuación con $n$ términos:
    \[
        c_1 e^{\lambda_1 t} + c_2 e^{\lambda_2 t} + \dots + c_n e^{\lambda_n t} = 0.
    \]
    Derivando ambos lados de la ecuación con respecto a $t$, obtenemos:
    \[
        c_1 \lambda_1 e^{\lambda_1 t} + c_2 \lambda_2 e^{\lambda_2 t} + \dots + c_n \lambda_n e^{\lambda_n t} = 0.
    \]
    Por otro lado, multiplicando la ecuación original por $\lambda_n$, tenemos:
    \[
        c_1 \lambda_n e^{\lambda_1 t} + c_2 \lambda_n e^{\lambda_2 t} + \dots + c_n \lambda_n e^{\lambda_n t} = 0.
    \]
    Restando esta última ecuación de la ecuación derivada, el término enésimo se cancela:
    \[
        c_1 (\lambda_1 - \lambda_n) e^{\lambda_1 t} + c_2 (\lambda_2 - \lambda_n) e^{\lambda_2 t} + \dots + c_{n-1} (\lambda_{n-1} - \lambda_n) e^{\lambda_{n-1} t} = 0.
    \]
    Esta es una combinación lineal de $n-1$ exponenciales que es igual a cero. Por nuestra hipótesis de inducción, los coeficientes que acompañan a cada exponencial deben ser nulos:
    \[
        c_j (\lambda_j - \lambda_n) = 0 \quad \text{para todo } j \in \{1, 2, \dots, n-1\}.
    \]
    Dado que los $\lambda_j$ son distintos dos a dos por hipótesis, sabemos que $\lambda_j - \lambda_n \neq 0$. Por lo tanto, podemos dividir por este factor y deducir que:
    \[
        c_j = 0 \quad \text{para todo } j \in \{1, 2, \dots, n-1\}.
    \]
    Sustituyendo estos resultados en la ecuación original de $n$ términos, nos queda simplemente $c_n e^{\lambda_n t} = 0$. Al igual que en el caso base, como $e^{\lambda_n t} \neq 0$, esto implica que $c_n = 0$. 
    
    Con esto quedan demostrados todos los coeficientes nulos, completando el paso inductivo y la demostración del lema.
\end{proof}

\subsection*{Sistemas Lineales Homogéneos: Soluciones Fundamental y Diagonalización}

Consideremos un sistema de ecuaciones diferenciales lineales homogéneo de primer orden con coeficientes constantes, dado por la expresión:
\[
    h' = Ah
\]
donde $A$ es una matriz de tamaño $n \times n$ con entradas reales o complejas, y $h(t)$ es un vector columna de $n$ funciones desconocidas. 

Supongamos que la matriz $A$ posee $n$ autovalores distintos dos a dos, denotados por $\lambda_1, \lambda_2, \dots, \lambda_n$, con sus respectivos autovectores asociados $v_1, v_2, \dots, v_n$. Sabemos que para cada par autovalor-autovector, la función vectorial $h_j(t) = e^{\lambda_j t} v_j$ es una solución del sistema.

A continuación, demostraremos que este conjunto de $n$ soluciones forma una base para el espacio de soluciones del sistema.

\subsubsection*{1. Independencia lineal de las soluciones}

Para verificar que las soluciones $\{ e^{\lambda_1 t} v_1, e^{\lambda_2 t} v_2, \dots, e^{\lambda_n t} v_n \}$ son linealmente independientes en todo $\mathbb{R}$, planteamos una combinación lineal igualada al vector nulo:
\[
    c_1 e^{\lambda_1 t} v_1 + c_2 e^{\lambda_2 t} v_2 + \dots + c_n e^{\lambda_n t} v_n = \mathbf{0} \quad \forall t \in \mathbb{R}
\]

Podemos agrupar los escalares $c_j$ con los vectores constantes $v_j$, reescribiendo la ecuación vectorial como:
\[
    e^{\lambda_1 t} (c_1 v_1) + e^{\lambda_2 t} (c_2 v_2) + \dots + e^{\lambda_n t} (c_n v_n) = \mathbf{0} \quad \forall t \in \mathbb{R}
\]

Dado que esta igualdad vectorial debe cumplirse componente a componente, podemos aplicar el lema previo (que establece la independencia lineal de funciones exponenciales con exponentes distintos) a cada una de las $n$ coordenadas del sistema. Al hacerlo, deducimos que los vectores coeficientes deben ser nulos:
\[
    c_j v_j = \mathbf{0} \quad \forall j \in \{1, 2, \dots, n\}
\]

Por definición, un autovector nunca es el vector nulo ($v_j \neq \mathbf{0}$). En consecuencia, la única forma de que se cumpla la igualdad anterior es que los escalares sean cero:
\[
    c_j = 0 \quad \forall j \in \{1, 2, \dots, n\}
\]

Esto demuestra que las $n$ soluciones son linealmente independientes. Si denotamos por $\Sigma_0$ al espacio de todas las soluciones del sistema homogéneo, concluimos que su dimensión es exactamente $n$, y que está generado por nuestro conjunto de soluciones:
\[
    \dim(\Sigma_0) = n, \quad \Sigma_0 = \mathcal{L}(\{e^{\lambda_j t} v_j : j = 1, 2, \dots, n\})
\]

\subsubsection*{2. Matriz Fundamental y el Wronskiano}

Dado que hemos encontrado un conjunto fundamental de soluciones, podemos construir la \textbf{matriz fundamental} del sistema, $\Phi(t)$, colocando estas soluciones como columnas de una matriz $n \times n$:
\[
    \Phi(t) = \begin{pmatrix} e^{\lambda_1 t} v_1 & e^{\lambda_2 t} v_2 & \dots & e^{\lambda_n t} v_n \end{pmatrix}
\]

Para garantizar que esta matriz es invertible (es decir, que forma un conjunto fundamental válido en todo instante $t$), calculamos su determinante (también conocido como el Wronskiano del sistema). Utilizando las propiedades del determinante para extraer escalares de las columnas, obtenemos:
\[
    |\Phi(t)| = e^{(\lambda_1 + \lambda_2 + \dots + \lambda_n) t} |v_1 v_2 \dots v_n|
\]

Dado que la función exponencial nunca se anula, y que los vectores $v_1, \dots, v_n$ son linealmente independientes (por estar asociados a autovalores distintos), su determinante $|v_1 v_2 \dots v_n|$ es distinto de cero. Por lo tanto:
\[
    |\Phi(t)| \neq 0 \quad \forall t \in \mathbb{R}
\]

\subsubsection*{3. Base de autovectores y Diagonalización}

Si evaluamos la matriz fundamental en el instante inicial $t = 0$, las exponenciales se evalúan a 1, y nos queda la matriz formada puramente por los autovectores:
\[
    \Phi(0) = \begin{pmatrix} v_1 & v_2 & \dots & v_n \end{pmatrix}
\]

Como demostramos que el determinante de esta matriz es distinto de cero, el conjunto de autovectores forma un sistema generador y libre para el espacio vectorial completo. Es decir:
\[
    \mathbb{C}^n = \mathcal{L} [v_1, v_2, \dots, v_n]
\]
\[
    \mathcal{B}_{\mathbb{C}} = \{v_1, v_2, \dots, v_n\} \text{ es una base de } \mathbb{C}^n
\]

Finalmente, si representamos la matriz original $A$ del sistema (la cual define una transformación lineal) expresada en esta nueva base $\mathcal{B}_{\mathbb{C}}$, obtenemos su forma diagonal, donde los elementos de la diagonal principal son precisamente los autovalores del sistema:
\[
    A_{\mathcal{B}_{\mathbb{C}}} = \begin{pmatrix} 
    \lambda_1 & 0 & \dots & 0 \\ 
    0 & \lambda_2 & \dots & 0 \\ 
    \vdots & \vdots & \ddots & \vdots \\ 
    0 & 0 & \dots & \lambda_n 
    \end{pmatrix}
\]

Este proceso evidencia cómo la estructura algebraica de la matriz (sus autovalores y autovectores) determina completamente la naturaleza geométrica y analítica de las soluciones del sistema diferencial.

\[
    A v_j = \lambda_j v_j \quad v_j \in \ker(A - \lambda_j I) \setminus \{0\} \quad \forall j \in \{1, 2, \dots, n\}
\]
\[
    \mathbb{C}^n = \bigoplus_{j=1}^n \ker(A - \lambda_j I) \quad \text{(descomposición espectral)}
\]
\[
    \dim(\ker(A - \lambda_j I) = 1 \quad \forall j \in \{1, 2, \dots, n\}
\]
\[
    u \in \mathbb{C}^n, \quad u = \sum_{j=1}^n c_j v_j \quad \text{(expresión de cualquier vector en la base de autovectores)}
\]

Ahora si $\sigma(A) = \{\lambda_1, \lambda_2, \dots, \lambda_q\}$ con $q < n$ autovalores distintos, cada uno con su respectiva multiplicidad algebraica $m_j$, tenemos:
\[
    \mathbb{C}^n = \bigoplus_{j=1}^q \ker((A - \lambda_j I)^{m_j}), \quad n = \sum_{j=1}^q m_j
\]
\[
    \ker(A - \lambda_j I) = \mathcal{L} [v_{j,1}, v_{j,2}, \dots, v_{j,m_j}]
\]
$e^{\lambda_j t} v_{j,l}$ es una solución del sistema homogéneo para cada $l \in \{1, 2, \dots, m_j\}$, y el conjunto de todas estas soluciones para $j = 1, 2, \dots, q$ forma una base del espacio de soluciones $\Sigma_0$.

\begin{teorema}
    Sea $A \in \mathcal{M}_{n \times n}(\mathbb{C})$ una matriz \textbf{diagonalizable} que tiene $q$ autovalores distintos $\lambda_1, \lambda_2, \dots, \lambda_q$ con multiplicidades algebraicas (y geométricas) $m_1, m_2, \dots, m_q$, de modo que $\sum_{j=1}^q m_j = n$. 
    
    Sean $\{v_{j,1}, v_{j,2}, \dots, v_{j,m_j}\}$ bases para cada uno de los subespacios propios asociados a $\lambda_j$. Entonces, el conjunto de $n$ funciones vectoriales:
    \[
        \mathcal{B} = \{ e^{\lambda_j t} v_{j,l} : j = 1, 2, \dots, q; \ l = 1, 2, \dots, m_j \}
    \]
    forma una base del espacio de soluciones $\Sigma_0$ del sistema lineal homogéneo $h' = Ah$.
\end{teorema}

\begin{proof}
    Para demostrar que $\mathcal{B}$ es una base de $\Sigma_0$, debemos probar dos cosas: que sus elementos son soluciones del sistema y que son linealmente independientes. Dado que la teoría general de sistemas lineales nos garantiza que la dimensión de $\Sigma_0$ es $n$, bastará con comprobar que estas $n$ funciones son soluciones y son linealmente independientes.

    \textbf{Verificación de las soluciones.} \\
    Para cualquier $j \in \{1, \dots, q\}$ y $l \in \{1, \dots, m_j\}$, calculemos la derivada de la función $h_{j,l}(t) = e^{\lambda_j t} v_{j,l}$:
    \[
        h_{j,l}'(t) = \lambda_j e^{\lambda_j t} v_{j,l}
    \]
    Por otro lado, como $v_{j,l}$ es un autovector asociado a $\lambda_j$, sabemos por definición que $A v_{j,l} = \lambda_j v_{j,l}$. Si evaluamos la acción de la matriz $A$ sobre nuestra función, obtenemos:
    \[
        A h_{j,l}(t) = A (e^{\lambda_j t} v_{j,l}) = e^{\lambda_j t} A v_{j,l} = e^{\lambda_j t} \lambda_j v_{j,l} = h_{j,l}'(t)
    \]
    Por lo tanto, cada función del conjunto $\mathcal{B}$ satisface la ecuación $h' = Ah$ y es, efectivamente, una solución del sistema homogéneo.

    \textbf{Independencia lineal.} \\
    Planteamos una combinación lineal de todos los elementos del conjunto igualada al vector nulo para todo instante $t \in \mathbb{R}$:
    \[
        \sum_{j=1}^q \sum_{l=1}^{m_j} c_{j,l} e^{\lambda_j t} v_{j,l} = \mathbf{0} \quad \forall t \in \mathbb{R}
    \]
    Podemos reorganizar esta suma agrupando los términos que comparten la misma función exponencial $e^{\lambda_j t}$:
    \[
        \sum_{j=1}^q \underbrace{\left( \sum_{l=1}^{m_j} c_{j,l} v_{j,l} \right)}_{u_j} e^{\lambda_j t} = \mathbf{0} \quad \forall t \in \mathbb{R}
    \]
    Llamemos $u_j$ al vector constante resultante de la suma interior. La ecuación vectorial se reescribe como:
    \[
        u_1 e^{\lambda_1 t} + u_2 e^{\lambda_2 t} + \dots + u_q e^{\lambda_q t} = \mathbf{0} \quad \forall t \in \mathbb{R}
    \]
    Por el lema previo (independencia lineal de funciones exponenciales con exponentes distintos), deducimos necesariamente que cada uno de los vectores coeficientes debe ser el vector nulo:
    \[
        u_j = \mathbf{0} \implies \sum_{l=1}^{m_j} c_{j,l} v_{j,l} = \mathbf{0} \quad \text{para cada } j \in \{1, 2, \dots, q\}
    \]
    
    \textbf{Conclusión sobre los coeficientes.} \\
    Fijemos un índice $j$. La suma $\sum_{l=1}^{m_j} c_{j,l} v_{j,l} = \mathbf{0}$ representa una combinación lineal nula de los autovectores asociados a $\lambda_j$. Por hipótesis, el conjunto $\{v_{j,1}, \dots, v_{j,m_j}\}$ es una base del subespacio propio asociado a $\lambda_j$ y, por tanto, son vectores linealmente independientes entre sí. Esto obliga a que todos sus coeficientes escalares sean nulos:
    \[
        c_{j,l} = 0 \quad \text{para todo } l \in \{1, 2, \dots, m_j\}
    \]
    Como este razonamiento es válido para los $q$ autovalores, concluimos que todos los coeficientes $c_{j,l}$ de la combinación lineal original son idénticamente cero.

    Por consiguiente, las $n$ soluciones propuestas son linealmente independientes y forman una base de $\Sigma_0$.
    
    \vspace{0.5cm}
    \textit{Nota de estructura:} En esta base formada íntegramente por autovectores, la matriz original $A$ adopta su forma diagonal (diagonal por bloques escalares), lo cual ilustra el desacoplamiento total del sistema:
    \[
        A_{\mathcal{B}} = \begin{pmatrix}
        \lambda_1 I_{m_1} & \mathbf{0} & \dots & \mathbf{0} \\
        \mathbf{0} & \lambda_2 I_{m_2} & \dots & \mathbf{0} \\
        \vdots & \vdots & \ddots & \vdots \\
        \mathbf{0} & \mathbf{0} & \dots & \lambda_q I_{m_q}
        \end{pmatrix}
    \]
\end{proof}
